\documentclass[11pt,a4paper]{article}
\usepackage[T1]{fontenc}
\usepackage{lmodern}
\usepackage[margin=25mm]{geometry}
\usepackage{amsmath,amssymb,amsthm,mathtools}
\usepackage{booktabs,array,tabularx}
\usepackage{enumitem}
\usepackage{xcolor}
\usepackage{microtype}
\usepackage{listings}
\usepackage{fancyhdr}
\usepackage[colorlinks=true,linkcolor=blue!45!black,citecolor=blue!45!black,urlcolor=blue!45!black]{hyperref}
\usepackage[capitalise,nameinlink]{cleveref}

\definecolor{codebg}{HTML}{F7F8FA}
\lstset{basicstyle=\ttfamily\footnotesize,backgroundcolor=\color{codebg},
  frame=single,rulecolor=\color{black!15},breaklines=true,columns=fullflexible,
  keepspaces=true,showstringspaces=false,numbers=left,numberstyle=\tiny\color{black!50},
  xleftmargin=12pt,framesep=4pt}

\pagestyle{fancy}
\fancyhf{}
\fancyhead[L]{\small\sffamily Harmonic branching mass and Mills-type prime-representing functions}
\fancyhead[R]{\small\sffamily KEIO}
\fancyfoot[C]{\thepage}
\renewcommand{\headrulewidth}{0.3pt}

\newtheorem{theorem}{Theorem}[section]
\newtheorem{proposition}[theorem]{Proposition}
\newtheorem{lemma}[theorem]{Lemma}
\newtheorem{corollary}[theorem]{Corollary}
\theoremstyle{definition}
\newtheorem{definition}[theorem]{Definition}
\newtheorem{problem}[theorem]{Problem}
\newtheorem{hypothesis}[theorem]{Hypothesis}
\newtheorem{remark}[theorem]{Remark}
\newtheorem{example}[theorem]{Example}

\newcommand{\N}{\mathbb{N}}
\newcommand{\Z}{\mathbb{Z}}
\newcommand{\R}{\mathbb{R}}
\newcommand{\PP}{\mathcal{P}}
\newcommand{\W}{\mathcal{W}}
\newcommand{\fl}[1]{\lfloor #1\rfloor}
\newcommand{\eps}{\varepsilon}
\newcommand{\meas}{\operatorname{meas}}
\newcommand{\Ch}{\operatorname{Ch}}
\newcommand{\VK}{\mathrm{VK}}

\title{\bfseries Harmonic Branching Mass and Mills-Type\\ Prime-Representing Functions below the Exponent $20/11$\\[6pt]
\large\normalfont A mass--compactness principle for prime trees and the exponent problem\\ of Mills, Kuipers, Niven and Wright}
\author{KEIO\\[2pt]\small research development and manuscript preparation with Claude (Anthropic)}
\date{8 November 2026}

\begin{document}
\maketitle

\begin{abstract}
For a real number $c>1$ let $\W(c)$ be the set of $A>1$ such that $\fl{A^{c^n}}$ is prime for every $n\ge 1$.
Mills proved $\W(3)\neq\emptyset$ in 1947, and the admissible exponents have since been extended by
Kuipers, Niven, Wright, Matomäki ($c\ge 2$), Islam ($c\ge 1.946$) and Järviniemi ($c\ge 20/11$).
The published arguments are \emph{local}: each new prime is selected inside the single short interval
prescribed by the previous prime, so they require power-saving bounds for the exceptional set of
short intervals that are poor in primes.
We formulate a \emph{harmonic branching problem} for trees of nested prime intervals and prove a
\emph{mass--compactness principle}: if, at every level, the harmonic mass $\sum 1/x$ carried by
defective nodes is at most half of the guaranteed lower bound for the harmonic mass of that level, then the
tree has an infinite branch. The principle only asks the exceptional set to beat the accumulated logarithmic losses,
which grow like $\exp(O(n^2))$, while the exceptional sets coming from zero-density estimates shrink
doubly exponentially along the tree. Combining it with a Selberg-type mean-square estimate with
Vinogradov--Korobov savings we prove, unconditionally, that $\W(c)$ is uncountable, unbounded,
Lebesgue-null and nowhere dense for every $c>6/5$ using only the zero-density estimates of Ingham and
Huxley, and for every $c>15/13$ using in addition the zero-density estimate of Guth and Maynard.
Under the density hypothesis the same holds for every $c>1$. The results hold for bounded variable
exponent sequences as well. The threshold $15/13$ is exactly $1/(1-\theta)$ with $\theta=2/15$ the
current almost-all exponent for the prime number theorem in short intervals.
We also record an explicit chain of $22$ proven primes for the exponent $5/4$ (up to $326$ digits).
The full conjectural range $1<c\le 15/13$ remains open unconditionally.
\end{abstract}

\noindent\textbf{Keywords:} Mills' constant; prime-representing functions; primes in almost all short intervals;
zero-density estimates; König's lemma; harmonic mass.\\
\noindent\textbf{MSC 2020:} 11N05, 11A41, 11M26.

\tableofcontents

\section{Introduction}

\subsection{The exponent problem}
Mills \cite{Mil47} showed that there is a real number $A$ such that $\fl{A^{3^n}}$ is prime for every
$n\ge1$. The exponent $3$ comes from Ingham's theorem that $[x,x+x^{5/8+\eps}]$ contains primes for large $x$.
For $c>1$ and, more generally, for a sequence $\mathbf c=(c_k)_{k\ge1}$ of reals $c_k>1$ with
$C_n=c_1c_2\cdots c_n$, put
\begin{gather*}
 \W(c)=\{A>1:\ \fl{A^{c^n}}\ \text{is prime for all } n\ge1\},\\
 \W(\mathbf c)=\{A>1:\ \fl{A^{C_n}}\ \text{is prime for all } n\ge1\}.
\end{gather*}
The question for which $c$ the set $\W(c)$ is nonempty has a long history, summarised in
\cref{tab:history}. Kuipers \cite{Kui50} treated integers $c\ge 3$, Niven \cite{Niv51} real $c>8/3$ and
Wright \cite{Wri54} $c\ge 925/348$; the theorem of Baker, Harman and Pintz on primes in intervals
$[x,x+x^{0.525}]$ gives $c\ge 2.106$; Matomäki \cite{Mat10} reached $c\ge2$ by working with almost all short
intervals; Islam \cite{Isl15} reached $c\ge 1.946$ and Järviniemi \cite{Jar22} reached $c\ge 20/11$.
Saito \cite{Sai21} proved that $\W(\mathbf c)$ has Hausdorff dimension $1$ for bounded sequences with
$c_k\ge2$. All exponents $c>1$ are expected to be admissible, since Cramér's model predicts primes in
every interval $[x,x+(\log x)^{3}]$.

\begin{table}[ht]
\centering\small
\caption{Admissible exponents for Mills-type prime-representing functions $\fl{A^{c^n}}$.}
\label{tab:history}
\begin{tabular}{@{}lll@{}}
\toprule
Source & admissible $c$ & input \\
\midrule
Mills 1947 \cite{Mil47} & $c=3$ & Ingham: primes in $[x,x+x^{5/8}]$\\
Kuipers 1950 \cite{Kui50} & integers $c\ge3$ & same\\
Niven 1951 \cite{Niv51} & $c>8/3$ & same\\
Wright 1954 \cite{Wri54} & $c\ge 925/348\approx2.658$ & general representing functions\\
via Baker--Harman--Pintz 2001 \cite{BHP01} & $c\ge 2.106$ & primes in $[x,x+x^{0.525}]$\\
Matomäki 2010 \cite{Mat10} & $c\ge 2$ & large gaps are rare \cite{Mat07}\\
Islam 2015 \cite{Isl15} & $c\ge 1.946$ & refined exceptional sets\\
Järviniemi 2022 \cite{Jar22} & $c\ge 20/11\approx1.818$ & sum of gaps $\ge x^{0.45}$ is $O(x^{0.63+\eps})$\\
\midrule
this paper, \cref{thm:main} & $c>6/5$ & Ingham--Huxley density, VK zero-free region\\
this paper, \cref{thm:main-GM} & $c>15/13\approx1.1538$ & adds Guth--Maynard density \cite{GM24}\\
this paper, \cref{thm:DH} & $c>1$ & density hypothesis\\
\bottomrule
\end{tabular}
\end{table}

\subsection{Main results}

\begin{theorem}[unconditional, classical inputs]\label{thm:main}
Let $\mathbf c=(c_k)_{k\ge1}$ satisfy $6/5<\inf_k c_k\le\sup_k c_k<\infty$. Then $\W(\mathbf c)$ is uncountable,
unbounded, of Lebesgue measure zero and nowhere dense. In particular, for every real $c>6/5$ there are
uncountably many $A>1$ such that $\fl{A^{c^n}}$ is prime for every $n\ge1$.
\end{theorem}

\begin{theorem}[unconditional, with the Guth--Maynard estimate]\label{thm:main-GM}
The conclusion of \cref{thm:main} holds for every sequence with $15/13<\inf_k c_k\le\sup_k c_k<\infty$.
In particular $\W(c)\neq\emptyset$ for every $c>15/13$.
\end{theorem}

\begin{theorem}[conditional]\label{thm:DH}
Assume the density hypothesis $N(\sigma,T)\ll_\eps T^{2(1-\sigma)+\eps}$ for $1/2\le\sigma\le1$.
Then the conclusion of \cref{thm:main} holds for every sequence with $1<\inf_k c_k\le\sup_k c_k<\infty$.
\end{theorem}

The exponents $6/5$ and $15/13$ are $1/(1-\theta_*)$ with $\theta_*=1/6$ (Huxley) and $\theta_*=2/15$
(Guth--Maynard), the known exponents for the prime number theorem in almost all short intervals
\cite{Hux72,GM24,expdb}. In the range $15/13<c<20/11$ the theorems are new; for $c\ge 20/11$ they
reprove known statements by a different method.

\subsection{The new principle}
Every proof in \cref{tab:history} is local. Given $p_n$, the next prime $p_{n+1}$ must be chosen in
$J(p_n)=[p_n^{c},(p_n+1)^c-1)$, an interval of length $\asymp X^{\theta}$ at height $X=p_n^c$, with
$\theta=1-1/c$. One needs a prime $p_{n+1}\in J(p_n)$ whose own interval $J(p_{n+1})$ is again rich in primes.
Since $J(p_n)$ contains only about $X^\theta/\log X$ primes, the exceptional set of bad intervals in the
whole dyadic range must be smaller than $X^{\theta}$ in count, i.e.\ one needs a \emph{power saving}.
This is what limits the method to $c\ge 20/11$ at present.

We replace the local selection by a \emph{global} count. We start from \emph{all} primes in a dyadic
interval $[P,2P]$, keep the whole tree of admissible continuations, and measure each level by the
harmonic mass $\sum_x 1/x$. Good nodes multiply their mass by about $1/\log x$; bad nodes are discarded.
The accumulated loss after $n$ levels is $\exp(O(n^2+n\log\log P))$, whereas the harmonic mass of
\emph{all} bad integers at level $n$ is $\exp(-\kappa(\log Y_n)^{1/3-o(1)})$ with $\log Y_n\asymp c^{n}\log P$,
which is doubly exponentially small in $n$. König's lemma then turns ``every level is nonempty'' into an
infinite branch. No single node is ever required to have a good child. The abstract form of this is the
mass--compactness principle, \cref{thm:mass}, which answers the harmonic branching problem, \cref{prob:H}.
The logical structure of the paper is therefore: a new problem (\cref{prob:H}), a new method solving it
(\cref{thm:mass,cor:split}), and its application, through the estimates of \cref{sec:aa}, to the classical
exponent problem (\cref{sec:proof}).

\subsection{Status of the claims}
\Cref{thm:main,thm:main-GM} are unconditional. They rely on published zero-density estimates
(Ingham \cite{Ing40}, Huxley \cite{Hux72}, Guth--Maynard \cite{GM24}, published in 2026), on the Vinogradov--Korobov zero-free
region \cite{Ivi85,For02} and on the explicit formula \cite{Dav00}; everything else is proved here.
\Cref{thm:DH} is conditional. The problem of admissibility for all $c\in(1,15/13]$ is \emph{not}
solved unconditionally; see \cref{sec:audit} for an itemised audit. The numerical data of
\cref{sec:numerics} are an illustration and are not used in any proof; \cref{app:proth} records a stopped
side computation.

\section{The harmonic branching problem and the mass--compactness principle}\label{sec:mass}

\begin{definition}[branching system]\label{def:branching}
A \emph{branching system} consists of a finite nonempty set $V_1\subset\Z_{\ge2}$ (the roots), for each
level $n\ge1$ a map $x\mapsto\Ch_n(x)$ assigning to every integer $x\ge2$ a finite set
$\Ch_n(x)\subset\Z_{\ge2}$ of children such that $\Ch_n(x)\cap\Ch_n(x')=\emptyset$ for $x\neq x'$, and for
each level a set $G_n\subset\Z_{\ge2}$ of \emph{good} integers. The levels are defined recursively by
\[
 V_{n+1}=\bigcup_{x\in V_n\cap G_n}\Ch_n(x).
\]
The weight of a finite set $S$ is $w(S)=\sum_{x\in S}1/x$.
\end{definition}

\begin{problem}[harmonic branching problem]\label{prob:H}
Find verifiable conditions on the \emph{local branching factors}
$\lambda_n=\inf_{x\in V_n\cap G_n} w(\Ch_n(x))\,x$ and on the \emph{defect masses}
$\beta_n\ge w(V_n\setminus G_n)$ which guarantee an infinite branch $x_1\in V_1$, $x_{k+1}\in\Ch_k(x_k)$,
$x_k\in G_k$ for all $k\ge1$, without any information on which individual nodes are good.
\end{problem}

\begin{theorem}[mass--compactness principle]\label{thm:mass}
Consider a branching system and positive numbers $\lambda_n$, nonnegative numbers $\beta_n$ such that
for every $n\ge1$:
\begin{enumerate}[label=\textup{(\roman*)}]
 \item $w(\Ch_n(x))\ge\lambda_n/x$ for every $x\in V_n\cap G_n$;
 \item $w(V_n\setminus G_n)\le\beta_n$;
 \item $\displaystyle \beta_n\le\frac12\,w(V_1)\prod_{k=1}^{n-1}\frac{\lambda_k}{2}$.
\end{enumerate}
Then $w(V_n)\ge w(V_1)\prod_{k<n}(\lambda_k/2)>0$ for every $n$, and there is an infinite branch
$(x_k)_{k\ge1}$ with $x_1\in V_1$, $x_{k+1}\in\Ch_k(x_k)$ and $x_k\in G_k$ for all $k$.
\end{theorem}

\begin{proof}
Put $H_n=w(V_n)$ and $M_n=w(V_1)\prod_{k<n}(\lambda_k/2)$. We show $H_n\ge M_n$ by induction; $H_1=M_1$.
If $H_n\ge M_n$, then (ii) and (iii) give $w(V_n\cap G_n)\ge H_n-\beta_n\ge H_n-M_n/2\ge H_n/2$.
Since the sets $\Ch_n(x)$ are pairwise disjoint, (i) yields
\[
 H_{n+1}=\sum_{x\in V_n\cap G_n}w(\Ch_n(x))\ge\lambda_n\,w(V_n\cap G_n)\ge \frac{\lambda_n}{2}H_n\ge M_{n+1}.
\]
Hence every $V_n$ is nonempty. Each $V_n$ is finite, because $V_1$ and all $\Ch_n(x)$ are finite.
Every $y\in V_{n+1}$ has exactly one parent $x\in V_n\cap G_n$ with $y\in\Ch_n(x)$, by disjointness.
Thus the $V_n$ form a locally finite forest with finitely many roots and with nodes at every depth.
By König's lemma (applied after adjoining a common root) there is an infinite path
$x_1,x_2,\dots$ with $x_{k+1}\in\Ch_k(x_k)$; a node of $V_k$ has children only if it lies in $G_k$,
so $x_k\in G_k$ for all $k$.
\end{proof}

\begin{remark}[why the principle is global]
Condition (iii) compares the defect mass at level $n$ with the \emph{accumulated} mass of the whole level.
It is irrelevant whether the bad nodes are concentrated below a particular node; this is the point
where local arguments need power savings and \cref{thm:mass} does not.
\end{remark}

\begin{corollary}[binary splitting]\label{cor:split}
Consider a branching system with numbers $\lambda_n,\beta_n$ satisfying \textup{(i)} and \textup{(ii)} of
\cref{thm:mass}, and put $m_n=\frac13\,w(V_1)\prod_{k=1}^{n-1}(\lambda_k/6)$. Assume that for every $n\ge1$
\[
 \text{(iii$'$)}\qquad \beta_n\le\tfrac12\,m_n\qquad\text{and}\qquad \max_{x\in V_n}\frac1x\le m_n .
\]
Then there are uncountably many infinite branches $x_1\in V_1$, $x_{k+1}\in\Ch_k(x_k)$, $x_k\in G_k$, and any two
of them differ at some level.
\end{corollary}

\begin{proof}
A finite set $S$ of integers with $\max_{x\in S}1/x\le w(S)/3$ can be split into two disjoint parts of weight
at least $w(S)/3$ each: order $S$ and cut after the shortest initial segment of weight $\ge w(S)/3$; that
segment has weight at most $2w(S)/3$. For $s=(s_1,s_2,\dots)\in\{0,1\}^{\N}$ put $U_1^s=V_1$, let $V_n^s$ be the part of
$U_n^s$ selected by $s_n$, and $U_{n+1}^s=\bigcup_{x\in V_n^s\cap G_n}\Ch_n(x)$. Clearly $U_n^s,V_n^s\subset V_n$.
We show $w(U_n^s)\ge3m_n$ and $w(V_n^s)\ge m_n$ by induction. For $n=1$, $w(U_1^s)=w(V_1)=3m_1$, the split is
possible because $\max 1/x\le m_1$, and $w(V_1^s)\ge m_1$. If $w(V_n^s)\ge m_n$, then by (ii) and (iii$'$)
$w(V_n^s\cap G_n)\ge m_n-\beta_n\ge m_n/2$, so by (i) and disjointness
$w(U_{n+1}^s)\ge\lambda_n m_n/2=3m_{n+1}$; the split is possible because $\max_{x\in V_{n+1}}1/x\le m_{n+1}$, and
$w(V_{n+1}^s)\ge m_{n+1}$. Hence every level of the tree of $s$ is nonempty, and König's lemma gives an infinite
branch $b(s)$ in it. If $s\ne s'$ first differ at index $n$, the $n$-th nodes of $b(s)$ and $b(s')$ lie in the
two disjoint parts of $U_n^s=U_n^{s'}$, so $b(s)\neq b(s')$.
\end{proof}

\section{Nested prime intervals}\label{sec:nested}

For $c>1$ and an integer $x\ge2$ put
\[
 J_c(x)=\bigl[x^c,\,(x+1)^c-1\bigr),\qquad \PP_c(x)=\PP\cap J_c(x),
\]
where $\PP$ is the set of primes. The intervals $J_c(x)$, $x\ge2$, are pairwise disjoint.

\begin{lemma}[nesting]\label{lem:nest}
Let $\mathbf c=(c_k)$ with $c_k>1$, and let $(p_n)_{n\ge1}$ be integers $p_n\ge2$ with
$p_{n+1}\in J_{c_{n+1}}(p_n)$ for all $n\ge1$. Then $A=\lim_{n\to\infty}p_n^{1/C_n}$ exists, $A>1$, and
$\fl{A^{C_n}}=p_n$ for every $n\ge1$.
\end{lemma}

\begin{proof}
Let $a_n=p_n^{1/C_n}$ and $b_n=(p_n+1)^{1/C_n}$. From $p_{n+1}\ge p_n^{c_{n+1}}$ we get $a_{n+1}\ge a_n$, and
from $p_{n+1}+1<(p_n+1)^{c_{n+1}}$ we get $b_{n+1}<b_n$. Hence $a_n\le A\le b_{n+1}<b_n$ for every $n$,
that is $p_n\le A^{C_n}<p_n+1$. Finally $A\ge a_1=p_1^{1/c_1}>1$.
\end{proof}

\begin{lemma}[interval length]\label{lem:length}
For $c>1$ and $x\ge2$ one has $c\,x^{c-1}-1\le |J_c(x)|\le c\,(x+1)^{c-1}$.
\end{lemma}

\begin{proof}
By the mean value theorem $(x+1)^c-x^c=c\,\xi^{c-1}$ with $x<\xi<x+1$, and $|J_c(x)|=(x+1)^c-x^c-1$.
\end{proof}

\section{Primes in almost all short intervals with Vinogradov--Korobov savings}\label{sec:aa}

Let $N(\sigma,T)$ denote the number of zeros $\rho=\beta+i\gamma$ of $\zeta(s)$ with $\beta\ge\sigma$ and
$|\gamma|\le T$. Write $\mathcal L(X)=(\log X)^{1/3}(\log\log X)^{-1/3}$.

\begin{hypothesis}[$\mathrm Z(\theta_1)$]\label{hyp:Z}
Let $0<\theta_1<1$. There are $\sigma_1\in[1/2,1)$, $\delta_0>0$ and $B\ge0$ such that, with
$a_*=2/(1-\theta_1)-\delta_0$,
\begin{enumerate}[label=\textup{(\alph*)}]
\item for every $\eps>0$: $N(\sigma,T)\ll_\eps T^{a_*(1-\sigma)+\eps}$ uniformly for $1/2\le\sigma\le\sigma_1$, $T\ge2$;
\item $N(\sigma,T)\ll T^{a_*(1-\sigma)}(\log T)^B$ uniformly for $\sigma_1\le\sigma\le1$, $T\ge2$.
\end{enumerate}
\end{hypothesis}

We use the following classical facts.
\begin{itemize}[leftmargin=1.5em]
\item[(F1)] Explicit formula \cite[Ch.~17]{Dav00}: for $x\ge2$ and $T\ge2$,
$\psi(x)=x-\sum_{|\gamma|<T}x^\rho/\rho+O\bigl(x(\log xT)^2/T+\log x\bigr)$.
\item[(F2)] Zero counting \cite[Ch.~15]{Dav00}: $N(t+1)-N(t)\ll\log(|t|+2)$ for the number $N(t)$ of zeros with $0<\gamma\le t$.
\item[(F3)] Vinogradov--Korobov \cite[Thm.~6.1]{Ivi85}, \cite{For02}: there is $c_0>0$ such that $\zeta(\sigma+it)\neq0$ for
$\sigma\ge1-c_0(\log|t|)^{-2/3}(\log\log|t|)^{-1/3}$, $|t|\ge3$.
\item[(F4)] Ingham \cite{Ing40}, \cite[Ch.~11]{Ivi85}: $N(\sigma,T)\ll T^{3(1-\sigma)/(2-\sigma)}(\log T)^5$ for $1/2\le\sigma\le1$.
\item[(F5)] Huxley \cite{Hux72}, \cite[Ch.~11]{Ivi85}: $N(\sigma,T)\ll T^{3(1-\sigma)/(3\sigma-1)}(\log T)^{44}$ for $3/4\le\sigma\le1$.
\item[(F6)] Guth--Maynard \cite{GM24}, in the form recorded in \cite[``Guth--Maynard bound'']{expdb}: for each fixed $\sigma\in[7/10,8/10]$
and every $\eps>0$, $N(\sigma,T)\ll_{\sigma,\eps} T^{15(1-\sigma)/(3+5\sigma)+\eps}$.
\end{itemize}

\begin{lemma}[verification of $\mathrm Z(\theta_1)$]\label{lem:Z}
\begin{enumerate}[label=\textup{(\roman*)}]
\item $\mathrm Z(\theta_1)$ holds for every $\theta_1>1/6$, using \textup{(F4)} and \textup{(F5)} only.
\item $\mathrm Z(\theta_1)$ holds for every $\theta_1>2/15$, using \textup{(F4)}, \textup{(F5)} and \textup{(F6)}.
\item Under the density hypothesis, $\mathrm Z(\theta_1)$ holds for every $\theta_1>0$.
\end{enumerate}
\end{lemma}

\begin{proof}
(i) On $[1/2,3/4]$ the exponent $3/(2-\sigma)$ of (F4) is at most $12/5$, and on $[3/4,1]$ the exponent
$3/(3\sigma-1)$ of (F5) is at most $12/5$. Hence $N(\sigma,T)\ll T^{12(1-\sigma)/5}(\log T)^{44}$ on $[1/2,1]$.
If $\theta_1>1/6$ then $12/5<2/(1-\theta_1)$; take $\sigma_1=1/2$ and $\delta_0=2/(1-\theta_1)-12/5$ in (b).

(ii) Let $\theta_1>2/15$, so that $30/13<2/(1-\theta_1)$. On $[1/2,7/10]$, (F4) gives exponent
$3/(2-\sigma)\le30/13$ with a logarithmic factor; on $[7/10,4/5]$, (F6) gives $15/(3+5\sigma)\le30/13$ with a
factor $T^\eps$; on $[4/5,1]$, (F5) gives $3/(3\sigma-1)\le 15/7<30/13$ with a logarithmic factor. Take
$\sigma_1=4/5$ and $\delta_0=\tfrac12\bigl(2/(1-\theta_1)-30/13\bigr)$, so that $a_*=30/13+\delta_0$. Then (b) holds on
$[4/5,1]$ directly from (F5) with $B=44$, because $15/7<a_*$; on $[1/2,7/10]$ the logarithmic factor in (F4) is
$\ll_\eps T^\eps$, which gives (a) there. For (a) on $[7/10,4/5]$, where (F6) is a pointwise statement, use that $N(\sigma,T)$ is nonincreasing in
$\sigma$: for a mesh $7/10=s_0<s_1<\dots<s_r=4/5$ with $s_{j+1}-s_j\le\delta_0/20$ and $\sigma\in[s_j,s_{j+1}]$,
$N(\sigma,T)\le N(s_j,T)\ll T^{\frac{30}{13}(1-s_j)+\eps}\le T^{\frac{30}{13}(1-\sigma)+\frac{30}{13}\cdot\frac{\delta_0}{20}+\eps}$, and the extra
exponent is absorbed into $\delta_0(1-\sigma)$ because $1-\sigma\ge1/5$ and $\frac{30}{13}\cdot\frac1{20}<\frac15$.

(iii) Let $\theta_1>0$. The density hypothesis gives (a) with exponent $2<2/(1-\theta_1)$ on $[1/2,1]$. For (b)
choose $\sigma_1<1$ with $3/(3\sigma_1-1)<2/(1-\theta_1)$, which is possible because $3/(3\sigma-1)\to3/2$ as
$\sigma\to1$, and use (F5) on $[\sigma_1,1]$. Take $\delta_0$ smaller than both gaps. If the density hypothesis is
only available pointwise in $\sigma$, the mesh argument of (ii) gives the required uniformity.
\end{proof}

\begin{proposition}[mean square with VK savings; classical, cf.\ \cite{SV77,Zac98,BP00,GT25}]\label{prop:ms}
Assume $\mathrm Z(\theta_1)$ and let $\theta_1\le\theta_2<1$. There are $\kappa>0$ and $X_0$ such that for
$X\ge X_0$ and $X^{\theta_1}\le h\le X^{\theta_2}$,
\[
 \int_X^{2X}\bigl|\psi(u+h)-\psi(u)-h\bigr|^2\,du\le h^2X\exp\bigl(-\kappa\,\mathcal L(X)\bigr).
\]
\end{proposition}

\begin{proof}
Let $T=X^{1-\theta_1}\exp(2\sqrt{\log X})$; then $2\le T\le X$ for large $X$. By (F1), for $u\in[X,2X]$,
\begin{gather*}
 \psi(u+h)-\psi(u)-h=-S(u)+O\Bigl(\frac{X\log^2X}{T}+\log X\Bigr),\\
 S(u)=\sum_{|\gamma|<T}\frac{(u+h)^\rho-u^\rho}{\rho}=\int_u^{u+h}F(v)\,dv,
\end{gather*}
with $F(v)=\sum_{|\gamma|<T}v^{\rho-1}$. The error term is at most $h\exp(-\sqrt{\log X})$ for large $X$, since
$X\log^2X/T=X^{\theta_1}\log^2X\,e^{-2\sqrt{\log X}}$ and $h\ge X^{\theta_1}$.

By Cauchy--Schwarz, $|S(u)|^2\le h\int_u^{u+h}|F(v)|^2dv$, and integrating over $u\in[X,2X]$ (each $v$ lies
in at most $h$ units of $u$-measure),
\[
 \int_X^{2X}|S(u)|^2du\le h^2\int_X^{3X}|F(v)|^2dv=h^2\sum_{\rho,\rho'}\int_X^{3X}v^{\rho+\overline{\rho'}-2}dv .
\]
For $\rho=\beta+i\gamma$ and $\rho'=\beta'+i\gamma'$ the inner integral is at most $2X^{\beta+\beta'-1}$ in
absolute value, and also at most $6X^{\beta+\beta'-1}/|\gamma-\gamma'|$ by integrating the oscillating factor.
Hence it is $\le12X^{\beta+\beta'-1}/(1+|\gamma-\gamma'|)$. With $X^{\beta+\beta'}\le(X^{2\beta}+X^{2\beta'})/2$
and symmetry,
\[
 \int_X^{3X}|F|^2\le 12\sum_{|\gamma|<T}X^{2\beta-1}\sum_{|\gamma'|<T}\frac1{1+|\gamma-\gamma'|}
 \ll(\log T)^2\sum_{|\gamma|<T}X^{2\beta-1},
\]
by (F2). Consequently
\[
 \int_X^{2X}|\psi(u+h)-\psi(u)-h|^2du\ll h^2X\Bigl((\log X)^2\Sigma+e^{-2\sqrt{\log X}}\Bigr),
 \qquad \Sigma=\sum_{|\gamma|<T}X^{2\beta-2}.
\]
Zeros with $\beta<1/2$ contribute at most $N(T)X^{-1}\ll T\log T/X\le X^{-\theta_1/2}$ to $\Sigma$.
For $\beta\ge1/2$, partial summation gives
\[
 \sum_{\beta\ge1/2}X^{2\beta-2}=X^{-1}N(\tfrac12,T)+2\log X\int_{1/2}^1X^{2\sigma-2}N(\sigma,T)\,d\sigma .
\]
By (F3), $N(\sigma,T)=0$ for $\sigma>1-\eta$ with $\eta=c_0(\log T)^{-2/3}(\log\log T)^{-1/3}$.
Write $a_*(1-\theta_1)=2-\delta_1$ with $\delta_1=\delta_0(1-\theta_1)>0$. Then
$T^{a_*(1-\sigma)}X^{2\sigma-2}=\exp\bigl(-(1-\sigma)[\delta_1\log X-2a_*\sqrt{\log X}]\bigr)
\le\exp(-(1-\sigma)\tfrac{\delta_1}{2}\log X)$ for large $X$.
On $[\sigma_1,1-\eta]$ hypothesis (b) bounds the integrand by $(\log X)^B\exp(-\tfrac{\delta_1}{2}\eta\log X)$.
On $[1/2,\sigma_1]$ hypothesis (a), with $\eps=(1-\sigma_1)\delta_1/4$, bounds it by $X^{-(1-\sigma_1)\delta_1/4}$.
Since $\log T\asymp\log X$, we have $\eta\log X\gg\mathcal L(X)$, and therefore
$\Sigma\ll(\log X)^{B+1}\exp(-c'\mathcal L(X))$ for some $c'>0$. The proposition follows with any
$\kappa<c'$.
\end{proof}

\begin{definition}[good and bad integers]\label{def:good}
Let $c>1$ and $\eps\in(0,1/2]$. An integer $x\ge2$ is \emph{$(c,\eps)$-good} if
\[
 \#\PP_c(x)\ge(1-\eps)\,\frac{|J_c(x)|}{\log\bigl((x+1)^c\bigr)},
\]
and \emph{$(c,\eps)$-bad} otherwise.
\end{definition}

\begin{corollary}[bad integers are sparse]\label{cor:bad}
Let $1<c_-\le c_+<\infty$, $\eps\in(0,1/2]$, and assume $\mathrm Z(\theta_1)$ for some
$\theta_1<1-1/c_-$. There are $\kappa'>0$ and $Y_0$ such that for $Y\ge Y_0$ and every $c\in[c_-,c_+]$,
\[
 \#\{x\in[Y,2Y):\ x\ \text{is }(c,\eps)\text{-bad}\}\le Y\exp\bigl(-\kappa'\mathcal L(Y)\bigr).
\]
\end{corollary}

\begin{proof}
Put $\theta=1-1/c\in[1-1/c_-,1-1/c_+]$ and $K=\lceil100c_+/\eps\rceil$. Split $[Y,2Y)$ into $K$ blocks
$[y_i,y_{i+1})$ with $y_i=Y(1+i/K)$. For $x$ in block $i$, \cref{lem:length} gives
$c\,y_i^{c-1}-1\le|J_c(x)|\le c\,y_i^{c-1}(1+2/K)^{c-1}$ for $Y$ large, hence
$c\,y_i^{c-1}\ge(1-\eps/50)|J_c(x)|$. Put $h_i=(1-\eps/4)(c\,y_i^{c-1}-1)$, so that
$(1-0.3\eps)|J_c(x)|\le h_i\le(1-\eps/4)|J_c(x)|$ for large $Y$; the upper bound holds because
$c\,y_i^{c-1}-1\le c\,x^{c-1}-1\le|J_c(x)|$.

Let $x$ be bad and $u\in U(x)=[x^c,x^c+\tfrac\eps4|J_c(x)|]$. Then $(u,u+h_i]\subset[x^c,(x+1)^c-1]$, so the primes in
$(u,u+h_i]$ are those of $\PP_c(x)$ plus possibly the endpoint $(x+1)^c-1$. Prime powers $p^k$ with
$k\ge2$ in an interval of length $h_i$ at height $\asymp Y^c$ number $O(Y^{c/2-1}+\log Y)$. Hence
\[
 \psi(u+h_i)-\psi(u)\le\log\bigl((x+1)^c\bigr)\bigl(\#\PP_c(x)+1\bigr)+O\bigl((Y^{c/2-1}+\log Y)\log Y\bigr)
 \le(1-\eps)|J_c(x)|+o(\eps h_i),
\]
and therefore $\psi(u+h_i)-\psi(u)-h_i\le-0.7\eps|J_c(x)|+o(\eps h_i)\le-\tfrac{\eps}{2}h_i$.
The sets $U(x)$ are pairwise disjoint and have measure $\tfrac\eps4|J_c(x)|\ge\tfrac\eps4h_i$; they lie in
$[Y^c,(2Y+1)^c]$, which is covered by at most $2c_++2$ dyadic intervals $[X,2X]$ with $Y^c\le X\le3^{c_+}Y^c$.
For each of them $X^{\theta_1}\le h_i\le X^{\theta_2}$ with $\theta_2=1-1/(2c_+)$, which exceeds every $\theta\le1-1/c_+$, once $Y$ is large; \cref{prop:ms} is applied with this fixed pair $(\theta_1,\theta_2)$.
By Chebyshev's inequality and \cref{prop:ms}, the measure of
$\{u\in[X,2X]:|\psi(u+h_i)-\psi(u)-h_i|\ge\tfrac\eps2h_i\}$ is at most $4\eps^{-2}X\exp(-\kappa\mathcal L(X))$.
Hence the number of bad $x$ in block $i$ is
\[
 \le\frac{4}{\eps h_i}\,(2c_++2)\,4\eps^{-2}\,3^{c_+}Y^c\exp(-\kappa\mathcal L(Y^c))
 \ll_{c_+,\eps}Y\exp(-\kappa\mathcal L(Y^c)),
\]
because $h_i\gg Y^{c-1}$. Summing over the $K$ blocks and using $\mathcal L(Y^c)\ge\mathcal L(Y)$ for large $Y$
gives the claim for any $\kappa'<\kappa$ and large $Y$.
\end{proof}

\section{Proof of the main theorems}\label{sec:proof}

Let $\mathbf c=(c_k)$ with $c_-=\inf c_k$ and $c_+=\sup c_k$, and let $\theta_*\in\{1/6,\,2/15,\,0\}$ be the
threshold corresponding to \cref{thm:main}, \cref{thm:main-GM} and \cref{thm:DH}, so that the assumption is
$c_->1/(1-\theta_*)$, i.e.\ $1-1/c_->\theta_*$. Fix $\theta_1$ with $\theta_*<\theta_1<1-1/c_-$; by \cref{lem:Z},
$\mathrm Z(\theta_1)$ holds (unconditionally in the first two cases). Fix $\eps=1/4$.

\paragraph{The branching system.}
Let $P$ be large. Put $V_1=\PP\cap[P,2P]$, $\Ch_n(x)=\PP_{c_{n+1}}(x)$ and
$G_n=\{x:\ x\text{ is }(c_{n+1},\eps)\text{-good}\}$. The child sets of one level are disjoint because the
intervals $J_{c_{n+1}}(x)$ are. Put $D_n=c_2c_3\cdots c_n$ ($D_1=1$), $M_1=2P$ and $M_{n+1}=(M_n+1)^{c_{n+1}}$.
By induction $V_n\subset[P^{D_n},M_n]$. Writing $\log M_n=D_na_n$ and using $\log(M_n+1)\le\log M_n+1/M_n$ we get
$a_{n+1}\le a_n+1/M_n\le a_n+P^{-D_n}$; since $D_n\ge c_-^{\,n-1}\ge1+(n-1)(c_--1)$, this gives
$a_n\le\log(2P)+2/P\le\log(3P)$ for large $P$. Hence
\[
 V_n\subset R_n=\bigl[P^{D_n},(3P)^{D_n}\bigr].
\]

\paragraph{Local branching factor.}
Let $x\in V_n\cap G_n$ and $c=c_{n+1}$. Every child $y$ satisfies $y<(x+1)^c$, and by goodness and
\cref{lem:length},
\[
 w(\Ch_n(x))\ge\frac{\#\PP_c(x)}{(x+1)^c}\ge\frac{(1-\eps)(c\,x^{c-1}-1)}{c\log(x+1)\,(x+1)^c}
 \ge\frac{1-2\eps}{x\log(x+1)}\ge\frac{\lambda_n}{x},
 \qquad \lambda_n=\frac{1}{2D_n\log(4P)},
\]
for $P$ large, since $\log(x+1)\le D_n\log(4P)$ on $R_n$.

\paragraph{Defect mass.}
By \cref{cor:bad}, applied to dyadic blocks $[Y,2Y)$ covering $R_n$ (at most $2D_n+2$ of them, all with
$Y\ge P^{D_n}$),
\[
 w(V_n\setminus G_n)\le\sum_{\text{blocks}}\frac1Y\,Y e^{-\kappa'\mathcal L(Y)}
 \le(2D_n+2)\exp\bigl(-\kappa'\mathcal L(P^{D_n})\bigr)=:\beta_n .
\]

\paragraph{Initial mass.}
By the prime number theorem, $w(V_1)\ge(\pi(2P)-\pi(P))/(2P)\ge1/(4\log P)$ for large $P$.

\paragraph{Verification of (iii$'$).}
We have $c_-^{\,n-1}\le D_n\le c_+^{\,n-1}$. Write $m=n-1$ and $\ell=\log P$. Then, with $m_n$ as in
\cref{cor:split},
\begin{gather*}
 \tfrac12 m_n=\tfrac16w(V_1)\prod_{k<n}\frac{\lambda_k}{6}\ge\frac{1}{24\ell}\,\bigl(12\log 4P\bigr)^{-m}\,c_+^{-m(m-1)/2},\\
 \beta_n\le(2c_+^{m}+2)\exp\Bigl(-\kappa'\frac{(c_-^{m}\ell)^{1/3}}{(\log(c_+^m\ell))^{1/3}}\Bigr).
\end{gather*}
Taking logarithms, it suffices that
\begin{multline*}
 \kappa'\frac{c_-^{m/3}\ell^{1/3}}{(m\log c_++\log\ell)^{1/3}}\\
 \ge\log(2c_+^m+2)+\log(24\ell)+m\log(12\log4P)+\tfrac{m(m-1)}2\log c_+
\end{multline*}
for all $m\ge0$. The right side is $\le C(1+m^2)(1+\log\ell)$ with $C=C(c_+)$, and the left side is
$\ge\kappa' c_-^{m/3}\ell^{1/3}(m+1)^{-1/3}(\log c_++\log\ell)^{-1/3}$. Since
$\mu=\inf_{m\ge0}c_-^{m/3}(m+1)^{-7/3}>0$, the inequality holds for all $m$ as soon as
$\kappa'\mu\,\ell^{1/3}\ge 2C(1+\log\ell)(\log c_++\log\ell)^{1/3}$, which is true for large $\ell$.
The second condition in (iii$'$) holds because $\max_{x\in V_n}1/x\le P^{-D_n}\le P^{-c_-^{m}}=e^{-c_-^m\ell}$,
whereas $m_n\ge e^{-C'(1+m^2)(1+\log\ell)}$; for large $\ell$ the first is smaller for every $m\ge0$.

\paragraph{Conclusion.}
\Cref{cor:split} gives uncountably many infinite branches $(p_n)$ with $p_{n+1}\in J_{c_{n+1}}(p_n)$, any two
of which differ at some level. By \cref{lem:nest} each branch yields $A\in\W(\mathbf c)$ with
$\fl{A^{C_n}}=p_n$, so distinct branches yield distinct $A$. Thus $\W(\mathbf c)$ is uncountable. Since
$A^{c_1}\in[P,2P+1)$ and $P$ can be taken arbitrarily large, $\W(\mathbf c)$ is unbounded.

\begin{remark}[size of the constants]
The constants $A$ produced here are large: $A^{c_1}\ge P$ with $P$ beyond the range where the estimates of
\cref{sec:aa} take effect. This is unavoidable for small $c$ in general: for $A<2$ the number $\fl{A^{c}}$ is tiny
and small cases can fail; for instance the comment to OEIS A243358 \cite{OEIS243358} records that for $A<2$ and
$c<1.2209864\ldots$ the sequence $\fl{A^{c^n}}$ must contain composite terms. No such restriction applies to large $A$.
The threshold $P_0$ is effective in principle (all analytic inputs have effective constants), but the
constant $A$ itself is obtained through König's lemma; unlike Mills' greedy construction, the proof gives no
algorithm with a known running time for computing an element of $\W(\mathbf c)$.
\end{remark}

\paragraph{Measure and category.}
Let $1<a<b$ and $W_n=\{A\in[a,b]:\ \fl{A^{C_k}}\in\PP\ (k\le n)\}$. Each $A\in W_n$ lies in an interval
$[p^{1/C_n},(p+1)^{1/C_n})$ with $p\in\PP$, $a^{C_n}\le p+1$, $p\le b^{C_n}$, whose length is
$\le p^{1/C_n}/(C_np)\le b/(C_np)$. Hence, by Mertens' theorem,
\[
 \meas W_n\le\frac{b}{C_n}\sum_{a^{C_n}-1\le p\le b^{C_n}}\frac1p\le\frac{b}{C_n}\Bigl(\log\frac{\log b}{\log a}+o(1)\Bigr)\longrightarrow0 .
\]
The closure of $W_n$ consists of $W_n$ and finitely many endpoints, so
$\overline{\W(\mathbf c)\cap[a,b]}\subset\bigcap_n\overline{W_n}$ has measure zero. A closed null set has empty
interior; hence $\W(\mathbf c)$ is Lebesgue-null and nowhere dense. This proves
\cref{thm:main,thm:main-GM,thm:DH}. \qed

\subsection{The weakest input: a transference theorem}
The proof above uses \cref{prop:ms} only through \cref{cor:bad}, and it needs far less than the
Vinogradov--Korobov saving. The next theorem isolates what is actually required: a lower bound of the right
order of magnitude (not an asymptotic) in almost all of the intervals $J_c(x)$, with an exceptional density
that beats every fixed power of $\log$ by a quasi-polynomial factor.

\begin{theorem}[transference]\label{thm:transfer}
Let $1<c_-\le c_+<\infty$ and $\delta\in(0,1]$. Suppose there are $C>\log c_+/(2(\log c_-)^2)$ and $Y_0$ such that
for all $Y\ge Y_0$ and all $c\in[c_-,c_+]$
\[
 \text{(S)}\qquad \#\Bigl\{x\in[Y,2Y):\ \#\PP_c(x)<\delta\,\frac{|J_c(x)|}{\log((x+1)^c)}\Bigr\}\le Y\exp\bigl(-C(\log\log Y)^2\bigr).
\]
Then for every sequence $\mathbf c$ with values in $[c_-,c_+]$ the set $\W(\mathbf c)$ is uncountable, unbounded,
Lebesgue-null and nowhere dense.
\end{theorem}

\begin{proof}
Run the construction of this section with $G_n$ the set of integers satisfying the inequality opposite to (S)
for $c=c_{n+1}$. The local branching factor becomes $\lambda_n=\delta/(2D_n\log 4P)$, and the defect mass is
$\beta_n\le(2D_n+2)\exp(-C(\log D_n+\log\ell)^2)$ with $\ell=\log P$. With $m=n-1$ and $\log D_n\ge m\log c_-$,
condition (iii$'$) of \cref{cor:split} follows from
\[
 \bigl[C\log^2c_--\tfrac12\log c_+\bigr]m^2+\bigl[2C\log c_-\log\ell-\log c_+-\log(24\ell/\delta)\bigr]m
 +\bigl[C\log^2\ell-2-\log(24\ell)\bigr]\ge0,
\]
which holds for all $m\ge0$ once $\ell$ is large, because each bracket is then nonnegative. The weight
condition in (iii$'$) and the remaining steps are as before.
\end{proof}

\begin{corollary}\label{cor:transfer}
\Cref{cor:bad} gives (S) with $\delta=\tfrac12$ and the much stronger bound $Y\exp(-\kappa'\mathcal L(Y))$ whenever
$c_->A^*/2$, where $A^*$ is a uniform zero-density exponent available with logarithmic losses near $\sigma=1$
(\cref{lem:Z}). This recovers \cref{thm:main,thm:main-GM,thm:DH}. Any future proof of (S) for some
$\theta=1-1/c$ below $2/15$, for instance a quantitative form of a sieve lower bound for primes in almost all
short intervals, would immediately extend the theorems to that $c$.
\end{corollary}

\begin{remark}[how far towards $c=1$]\label{rem:towards1}
For constant $c$ the condition on $C$ reads $C>1/(2\log c)$, and the exceptional density needed is
$\exp(-C(\log\log Y)^2)=(\log Y)^{-C\log\log Y}$: super-polylogarithmic, but far weaker than the
Vinogradov--Korobov saving. As $c\downarrow1$ the constant $C$ grows like $1/(2(c-1))$, but the saving stays of
quasi-polylogarithmic type; what becomes hard is the interval length $X^{1-1/c}$. The obstruction to $c\to1$ is therefore entirely the lack of almost-all prime
information for intervals $[x,x+x^\theta]$ with small $\theta$: lower bounds of the right order are known in almost
all such intervals only for $\theta>1/22$ \cite{Li25b}, and asymptotics only for $\theta>2/15$. Reaching every
$c>1$ by this method requires (S) for every $\theta>0$, which follows from the density hypothesis and is
otherwise open.
\end{remark}

\section{Local versus global: where the exponent comes from}\label{sec:compare}

\begin{table}[ht]
\centering\small
\caption{What each method needs from the exceptional set $E(y)$ of integers $x\in[y,2y)$ whose interval
$J_c(x)$ (length $\asymp y^{c-1}=(y^c)^{\theta}$, $\theta=1-1/c$) is poor in primes.}
\label{tab:compare}
\begin{tabularx}{\linewidth}{@{}lXl@{}}
\toprule
Method & requirement & present range\\
\midrule
local selection \cite{Mat10,Jar22} & $\#E(y)\le y^{\theta-\delta}$, i.e.\ fewer bad integers than primes in one interval $J_c(p)$, $p\asymp y^{1/c}$ & $c\ge20/11$\\
global mass (\cref{thm:mass}) & $\#E(y)\le y\exp(-C(\log\log y)^2)$ for a suitable $C$; VK savings suffice & $c>15/13$\\
\bottomrule
\end{tabularx}
\end{table}

The local method must find a good child inside one interval $J_c(p_n)$; if $p_{n+1}\asymp y$, this interval
contains about $y^{\theta}/\log y$ primes, while the bad integers are counted in the whole range $[y,2y)$.
It therefore needs a power saving of size $y^{1-\theta}=y^{1/c}$ in the exceptional set. Results of this
type come from large-value estimates for Dirichlet polynomials and from sieve methods; with the
bound of Järviniemi \cite{Jar22} they reach $\theta=0.45$, i.e.\ $c=20/11$.

The global method compares the bad mass with the mass of the whole level, which is only
$\exp(O(n^2+n\log\log P))$ smaller than the initial mass. The bad mass at level $n$ is
$\exp(-\kappa(\log Y_n)^{1/3-o(1)})$ with $\log Y_n\ge c_-^{n-1}\log P$. Any saving of the form
$\exp(-(\log Y)^{\delta})$ with fixed $\delta>0$ would suffice; polylogarithmic savings would not, because the accumulated loss
$\prod_{k<n}\log Y_k$ grows faster than every fixed power of $\log Y_n$ (for constant exponent $c$ it equals
$(\log Y_n)^{(1/2+o(1))\log\log Y_n/\log c}$ as $n\to\infty$). This explains the thresholds:
\[
 c>\frac{1}{1-\theta_{\mathrm{PNT}\text{-}\mathrm{AA}}},\qquad \theta_{\mathrm{PNT}\text{-}\mathrm{AA}}=\tfrac16\ \text{(Huxley)},\ \tfrac2{15}\ \text{(Guth--Maynard)},\ 0\ \text{(density hypothesis)}.
\]
Equivalently, writing $A^*=\sup_{1/2\le\sigma\le1}A(\sigma)$ for the best uniform zero-density exponent,
the zero-density route gives exactly
\[
 \W(c)\neq\emptyset\qquad\text{for every } c>A^*/2 ,
\]
since $1/(1-(1-2/A^*))=A^*/2$: $A^*=12/5$ gives $6/5$, $A^*=30/13$ gives $15/13$, and $A^*=2$ (the density
hypothesis) gives every $c>1$. \Cref{thm:transfer} below isolates the weakest input the method needs.
Below $c=15/13$ the only available almost-all results for intervals $[x,x+x^\theta]$, $\theta\le2/15$, are sieve
lower bounds, down to $\theta>1/22$ \cite{Li24,Li25b}. Their exceptional sets have not, to our knowledge, been
quantified beyond $o(X)$. If they satisfy hypothesis (S) of \cref{thm:transfer}, the method gives every
$c>22/21\approx1.0476$, but not $c\to1$: for $1-1/c\le1/22$ not even the existence of a prime in almost all
intervals $[x,x+x^{1-1/c}]$ is known.

\begin{remark}[relation to Mills' constant]
The theorems concern existence and do not identify any specific constant. In particular they say nothing
about the least element of $\W(3)$ (Mills' constant), whose irrationality was proved by Saito
\cite{Sai25} and whose transcendence is open \cite{Sai25b}.
\end{remark}

\section{A certified chain for the exponent 5/4}\label{sec:numerics}

The exponent $c=5/4$ lies in the new range $(15/13,20/11)$; it is covered by \cref{thm:main} since $5/4>6/5$.
For $c=5/4$ the interval condition is purely integral:
\[
 m\in J_{5/4}(x)\iff m^4\ge x^5\ \text{and}\ (m+1)^4<(x+1)^5 .
\]
Starting from $p_1=1009$ and taking each time the least prime of $J_{5/4}(p_n)$ gives
$22$ primes $p_1,\dots,p_{22}$, each proven prime by \texttt{isprime} in PARI/GP~2.15 \cite{PARI}, with all
nesting conditions checked in exact integer arithmetic (\cref{lst:gp}). The first terms are
\[
 1009,\ 5689,\ 49409,\ 736657,\ 21581503,\ 1470964273,\ 288073066607,\ 211046580457399,\ \dots
\]
and the digit counts are $4,4,5,6,8,10,12,15,18,23,28,35,44,55,69,86,107,134,167,209,261,326$.
The offsets $p_{n+1}-\lceil p_n^{5/4}\rceil$ never exceed $954$, while $|J_{5/4}(p_n)|\approx1.25\,p_n^{1/4}$;
for instance $|J_{5/4}(p_{21})|\approx 10^{65}$. Any $A$ with $\fl{A^{(5/4)^n}}=p_n$ for $n\le22$ satisfies
\[
 A=253.06488747154176508840743158311366901154901154046745858478949072882087\ldots
\]
(to the precision fixed by $p_{22}$). This finite computation illustrates the nested-interval mechanism; it is
not a proof that this greedy sequence continues forever, and \cref{thm:main} does not assert that.

\appendix
\section{A stopped side computation: Proth candidates with 30{,}103 digits}\label{app:proth}

Independently of the theorems, we searched for a prime larger than the $18{,}886$-digit Pocklington prime of the
author's earlier computations, in the form $N_k=k\cdot2^{100000}+1$ with odd $k\le2\cdot10^6$, where Proth's theorem
certifies primality by a single congruence $a^{(N_k-1)/2}\equiv-1\pmod{N_k}$. Sieving by all primes up to $10^{10}$
(\cref{lst:sieve}) left $48{,}707$ candidates. The test of \cref{lst:c} reduces modulo $k2^{n}+1$ without division
(writing $t=H2^{n}+L$ and $Q=\lfloor H/k\rfloor$, one has $t\equiv t-Q N_k$) and was validated against
\texttt{gmpy2} on smaller exponents. Each test took about $17$\,s on one core.

The run was stopped by the time limit of the computing environment after $1{,}639$ candidates
($k\le68{,}751$), all composite; no prime was found. Heuristically a sieved candidate is prime with probability
about $e^{\gamma}\log(10^{10})/\log N_k\approx5.9\cdot10^{-4}$, so the expected number of tests before the first
prime is about $1{,}700$, and the observed outcome has probability about $e^{-0.97}\approx0.38$. The logs of the
tested candidates and the list of the $47{,}068$ untested ones are distributed with the source files, so that the
search can be resumed without repetition. No claim is made on the basis of this computation.

\section{Audit against the author's success criteria}\label{sec:audit}

The criteria below were fixed in advance by the author. Each verdict is stated as strictly as we can.

\begin{table}[ht]
\centering\small
\begin{tabularx}{\linewidth}{@{}lXl@{}}
\toprule
Criterion & Assessment & Verdict\\
\midrule
G1 prior named problem & The exponent problem for Mills-type prime-representing functions, posed implicitly by Mills (1947) and
studied explicitly by Kuipers, Niven, Wright, Matomäki, Islam and Järviniemi; the statement
``$\exists A$: $\fl{A^{c^n}}$ prime $\forall n$'' predates this work. & met\\
G2 problem unchanged & For each fixed $c$ we prove the original statement, without modified constants or reformulation. & met\\
G3 actually open & For $c<20/11$ no proof was available in the literature we could find
(\cite{Jar22} and the exponent database \cite{expdb} list $20/11$ and $2/15$ as current records). & met for $c\in(15/13,20/11)$\\
G4 unconditional and complete & \Cref{thm:main,thm:main-GM}: complete unconditional proofs, all quantifiers closed, for every
$c>6/5$, resp.\ $c>15/13$. The full conjecture (all $c>1$) is proved only under the density hypothesis. &
met on $(15/13,\infty)$; \textbf{not} met on $(1,15/13]$\\
G5 not a restatement & The operative new ingredient is a quantitative non-vanishing estimate for the harmonic branching mass
(\cref{thm:mass}) combined with VK-saving exceptional-set bounds; no identity or change of variables carries the proof. & met\\
G6 no prior identical solution & We found no earlier global (mass--compactness) argument for prime-representing functions and no result
below $20/11$; absence in our search does not prove priority. & no counterexample found\\
\bottomrule
\end{tabularx}
\end{table}

\noindent In plain terms: the paper resolves the Mills exponent problem on the new range $15/13<c<20/11$, shows
(\cref{thm:transfer}, \cref{rem:towards1}) exactly which almost-all input would be needed to go further, and
does not resolve the problem for $1<c\le15/13$. With present knowledge no almost-all method can go below
$c=22/21$.

\section{Verification code}\label{app:code}

\begin{lstlisting}[language={},caption={PARI/GP script for the chain of \cref{sec:numerics}.},label={lst:gp}]
\\ Greedy chain for c = 5/4.  J(x) = { m : m^4 >= x^5 and (m+1)^4 < (x+1)^5 }.
\\ Each term is proven prime by isprime() (PARI/GP primality proof).
default(parisize, 10^8);
{
p = 1009; L = List([p]);
for (n = 1, 21,
  lo = sqrtnint(p^5, 4); if (lo^4 < p^5, lo++);
  m = nextprime(lo);
  if (!((m+1)^4 < (p+1)^5), error("stuck at n=", n));
  if (!isprime(m), error("not proven prime"));
  listput(L, m); p = m);
}
{
for (i = 1, #L,
  q = L[i];
  if (i == 1, off = 0, a = L[i-1]; r = sqrtnint(a^5, 4); if (r^4 < a^5, r++); off = q - r);
  printf("%d  digits=%d  offset=%d\n", i, #digits(q), off));
}
{
ok = 1;
for (i = 1, #L-1, a = L[i]; b = L[i+1]; if (!(b^4 >= a^5 && (b+1)^4 < (a+1)^5), ok = 0));
print("nesting ok: ", ok);
}
write("chain54.txt", Vec(L));
\p 400
{
N = #L; A = exp(log(L[N] + 1/2) * (4/5)^N);
print("A approx = ", A);
bad = 0; for (i = 1, N, if (floor(A^((5/4)^i)) != L[i], bad++));
print("floor mismatches: ", bad);
}
quit;
\end{lstlisting}

\begin{lstlisting}[language=Python,caption={Independent Proth certificate check (Python 3 with gmpy2), for use if the search is resumed.},label={lst:py}]
#!/usr/bin/env python3
"""Independent check of a Proth certificate N = k*2^n + 1 (k odd, k < 2^n).

Proth's theorem: if a^((N-1)/2) == -1 (mod N) for some integer a, then N is prime.
Usage: python3 verify_proth.py k n a
"""
import sys, hashlib, time
import gmpy2
from gmpy2 import mpz

def main():
    k, n, a = (int(x) for x in sys.argv[1:4])
    assert k % 2 == 1 and k > 0, "k must be odd and positive"
    assert k < 2**n, "Proth's theorem needs k < 2^n"
    N = mpz(k) * (mpz(1) << n) + 1
    t0 = time.time()
    r = gmpy2.powmod(a, (N - 1) // 2, N)
    dt = time.time() - t0
    digits = len(str(N))
    print(f"N = {k}*2^{n}+1 has {digits} decimal digits and {N.bit_length()} bits")
    print(f"Jacobi({a}, N) = {gmpy2.jacobi(a, N)}")
    print(f"a^((N-1)/2) mod N == N-1 : {r == N - 1}   ({dt:.1f} s)")
    print("SHA-256 of decimal expansion:", hashlib.sha256(str(N).encode()).hexdigest())
    if r == N - 1:
        print("PROVEN PRIME by Proth's theorem.")
    else:
        print("certificate FAILED")
        sys.exit(1)

if __name__ == "__main__":
    main()
\end{lstlisting}

\begin{lstlisting}[language=C,caption={Search program used to find the Proth prime (GMP; reduction modulo $k2^n+1$ without division).},label={lst:c}]
// Proth test for N = k*2^n+1 with fast reduction. Usage: proth n listfile slice nslices outfile [maxsteps]
#include <gmp.h>
#include <stdio.h>
#include <stdlib.h>
#include <time.h>
static void redmod(mpz_t t, const mpz_t N, unsigned long k, unsigned long n, mpz_t Q, mpz_t u){
  mpz_tdiv_q_2exp(Q, t, n);
  mpz_tdiv_q_ui(Q, Q, k);
  mpz_mul_ui(u, Q, k);
  mpz_mul_2exp(u, u, n);
  mpz_add(u, u, Q);
  mpz_sub(t, t, u);
  while (mpz_sgn(t) < 0) mpz_add(t, t, N);
  while (mpz_cmp(t, N) >= 0) mpz_sub(t, t, N);
}
// returns 1 if a^((N-1)/2) == -1, 0 if == +1 (inconclusive), -1 otherwise (composite)
static int proth_test(unsigned long k, unsigned long n, unsigned long a, long maxsteps){
  mpz_t N,x,t,Q,u; mpz_inits(N,x,t,Q,u,NULL);
  mpz_set_ui(N,k); mpz_mul_2exp(N,N,n); mpz_add_ui(N,N,1);
  mpz_set_ui(x,a); mpz_powm_ui(x,x,k,N);        // x = a^k mod N
  long steps = (maxsteps>0 && maxsteps < (long)n-1) ? maxsteps : (long)n-1;
  for (long i=0;i<steps;i++){ mpz_mul(t,x,x); redmod(t,N,k,n,Q,u); mpz_swap(x,t); }
  int res;
  if (maxsteps>0 && steps < (long)n-1) { res = 2; gmp_printf("%Zx\n", x); }
  else { mpz_add_ui(t,x,1); if (mpz_cmp(t,N)==0) res=1; else if (mpz_cmp_ui(x,1)==0) res=0; else res=-1; }
  mpz_clears(N,x,t,Q,u,NULL);
  return res;
}
static unsigned long pick_base(unsigned long k, unsigned long n){
  mpz_t N; mpz_init(N); mpz_set_ui(N,k); mpz_mul_2exp(N,N,n); mpz_add_ui(N,N,1);
  unsigned long ps[]={3,5,7,11,13,17,19,23,29,31,37,41,43,47,53,59,61,67,71,73,79,83,89,97};
  unsigned long a=0;
  for (int i=0;i<24;i++){ if (mpz_ui_kronecker(ps[i],N)==-1){ a=ps[i]; break; } }
  mpz_clear(N); return a;
}
int main(int argc,char**argv){
  unsigned long n=strtoul(argv[1],0,10);
  FILE*f=fopen(argv[2],"r"); int slice=atoi(argv[3]), ns=atoi(argv[4]);
  FILE*o=fopen(argv[5],"a"); long maxsteps = argc>6? atol(argv[6]) : 0;
  unsigned long k; long idx=0;
  while (fscanf(f,"%lu",&k)==1){
    if ((idx++ % ns)!=slice) continue;
    unsigned long a=pick_base(k,n);
    if (!a){ fprintf(o,"k=%lu nobase\n",k); fflush(o); continue; }
    time_t t0=time(0);
    int r=proth_test(k,n,a,maxsteps);
    fprintf(o,"k=%lu a=%lu res=%d secs=%ld\n",k,a,r,(long)(time(0)-t0)); fflush(o);
    if (r==1){ fprintf(o,"PRIME k=%lu n=%lu base=%lu\n",k,n,a); fflush(o); break; }
    if (maxsteps>0) break;
  }
  return 0;
}
\end{lstlisting}

\begin{lstlisting}[language=C,caption={Sieve used to select the Proth candidates.},label={lst:sieve}]
// Sieve odd k in [1,K] for N=k*2^n+1 against primes l in [3,B]. Output surviving k.
#include <stdio.h>
#include <stdlib.h>
#include <stdint.h>
#include <string.h>
#include <math.h>
static uint64_t mulmod(uint64_t a,uint64_t b,uint64_t m){return (unsigned __int128)a*b%m;}
static uint64_t powmod(uint64_t a,uint64_t e,uint64_t m){uint64_t r=1%m;a%=m;while(e){if(e&1)r=mulmod(r,a,m);a=mulmod(a,a,m);e>>=1;}return r;}
static int64_t inv(int64_t a,int64_t m){int64_t g=m,x=0,x1=1,a1=a;while(a1){int64_t q=g/a1,t=g-q*a1;g=a1;a1=t;t=x-q*x1;x=x1;x1=t;}if(x<0)x+=m;return x;}
int main(int argc,char**argv){
  uint64_t n=strtoull(argv[1],0,10), K=strtoull(argv[2],0,10), B=strtoull(argv[3],0,10);
  uint8_t *alive=malloc(K+1); memset(alive,1,K+1);
  for(uint64_t k=0;k<=K;k+=2) alive[k]=0; // even k
  // segmented sieve of Eratosthenes for primes up to B
  uint64_t sq=(uint64_t)sqrt((double)B)+1;
  uint8_t *small=calloc(sq+1,1); uint64_t *sp=malloc(sizeof(uint64_t)*sq); size_t ns=0;
  for(uint64_t i=2;i<=sq;i++){ if(!small[i]){ sp[ns++]=i; for(uint64_t j=i*i;j<=sq;j+=i) small[j]=1; } }
  const uint64_t SEG=1<<24; uint8_t *seg=malloc(SEG);
  uint64_t cnt=0;
  for(uint64_t lo=3; lo<=B; lo+=SEG){
    uint64_t hi=lo+SEG-1; if(hi>B) hi=B;
    memset(seg,0,SEG);
    for(size_t t=0;t<ns;t++){ uint64_t p=sp[t]; if(p*p>hi) break; uint64_t s=((lo+p-1)/p)*p; if(s<p*p) s=p*p; for(uint64_t j=s;j<=hi;j+=p) seg[j-lo]=1; }
    for(uint64_t l=lo;l<=hi;l++){ if(seg[l-lo]) continue; if(l==2) continue;
      // N = k*2^n+1 = 0 mod l  <=> k = -inv(2^n) mod l
      uint64_t t2=powmod(2,n,l); uint64_t r=(l-(uint64_t)inv((int64_t)t2,(int64_t)l))%l;
      for(uint64_t k=r;k<=K;k+=l){ if(k==0) continue; alive[k]=0; }
      cnt++;
    }
  }
  uint64_t surv=0; for(uint64_t k=1;k<=K;k++) if(alive[k]){ printf("%llu\n",(unsigned long long)k); surv++; }
  fprintf(stderr,"primes sieved %llu survivors %llu\n",(unsigned long long)cnt,(unsigned long long)surv);
  return 0;
}
\end{lstlisting}

\begin{thebibliography}{99}\small
\bibitem{BHP01} R.~C. Baker, G.~Harman, J.~Pintz, The difference between consecutive primes, II, \emph{Proc. London Math. Soc.} (3) 83 (2001), 532--562.
\bibitem{Dav00} H.~Davenport, \emph{Multiplicative Number Theory}, 3rd ed., revised by H.~L. Montgomery, Graduate Texts in Math. 74, Springer, 2000.
\bibitem{For02} K.~Ford, Vinogradov's integral and bounds for the Riemann zeta function, \emph{Proc. London Math. Soc.} (3) 85 (2002), 565--633.
\bibitem{GM24} L.~Guth, J.~Maynard, New large value estimates for Dirichlet polynomials, \emph{Ann. of Math.} (2) 203 (2026), 623--675; arXiv:2405.20552.
\bibitem{Hux72} M.~N. Huxley, On the difference between consecutive primes, \emph{Invent. Math.} 15 (1972), 164--170.
\bibitem{Ing40} A.~E. Ingham, On the estimation of $N(\sigma,T)$, \emph{Quart. J. Math. Oxford} 11 (1940), 201--202.
\bibitem{Isl15} A.~Islam, \emph{Multiplicative properties of integers in short intervals}, PhD thesis, Royal Holloway, University of London, 2015.
\bibitem{Ivi85} A.~Ivi\'c, \emph{The Riemann Zeta-Function}, Wiley, 1985.
\bibitem{Jar22} O.~Järviniemi, On large differences between consecutive primes, preprint, arXiv:2212.10965 (2022).
\bibitem{Kui50} L.~Kuipers, Prime-representing functions, \emph{Nederl. Akad. Wetensch. Proc.} 53 (1950), 309--310 $=$ \emph{Indag. Math.} 12 (1950), 57--58.
\bibitem{Li24} R.~Li, Primes in almost all short intervals, preprint, arXiv:2407.05651 (2024).
\bibitem{Li25b} R.~Li, Primes in almost all short intervals II (2025); exponent $1/22$ as recorded in \cite{expdb}.
\bibitem{SV77} B.~Saffari, R.~C. Vaughan, On the fractional parts of $x/n$ and related sequences II, \emph{Ann. Inst. Fourier} 27 (1977), 1--30.
\bibitem{Zac98} A.~Zaccagnini, Primes in almost all short intervals, \emph{Acta Arith.} 84 (1998), 225--244.
\bibitem{BP00} D.~Bazzanella, A.~Perelli, The exceptional set for the number of primes in short intervals, \emph{J. Number Theory} 80 (2000), 109--124.
\bibitem{GT25} A.~Gafni, T.~Tao, On the number of exceptional intervals to the prime number theorem in short intervals, preprint, arXiv:2505.24017 (2025).
\bibitem{Mat07} K.~Matomäki, Large differences between consecutive primes, \emph{Q. J. Math.} 58 (2007), 489--518.
\bibitem{Mat10} K.~Matomäki, Prime-representing functions, \emph{Acta Math. Hungar.} 128 (2010), 307--314.
\bibitem{Mil47} W.~H. Mills, A prime-representing function, \emph{Bull. Amer. Math. Soc.} 53 (1947), 604.
\bibitem{Niv51} I.~Niven, Functions which represent prime numbers, \emph{Proc. Amer. Math. Soc.} 2 (1951), 753--755.
\bibitem{OEIS243358} OEIS Foundation, Sequence A243358 (A.~V. Kulsha, D.~J. Broadhurst), \emph{The On-Line Encyclopedia of Integer Sequences}, \url{https://oeis.org/A243358}.
\bibitem{PARI} The PARI~Group, \emph{PARI/GP version 2.15.4}, Univ. Bordeaux, 2023, \url{https://pari.math.u-bordeaux.fr/}.
\bibitem{Sai21} K.~Saito, Prime-representing functions and Hausdorff dimension, \emph{Acta Math. Hungar.} (2021); arXiv:2102.04038.
\bibitem{Sai25} K.~Saito, Mills' constant is irrational, \emph{Mathematika} 71 (2025), e70027; arXiv:2404.19461.
\bibitem{Sai25b} K.~Saito, Transcendency of variants of Mills' constant, preprint, arXiv:2508.16068 (2025).
\bibitem{expdb} T.~Tao, T.~Trudgian, A.~Yang et al., \emph{Analytic Number Theory Exponent Database}, \url{https://github.com/teorth/expdb} (accessed 8 November 2026), chapters ``Zero density theorems'' and ``Distribution of primes: short ranges''.
\bibitem{Wri54} E.~M. Wright, A class of representing functions, \emph{J. London Math. Soc.} 29 (1954), 63--71.
\end{thebibliography}

\end{document}
